\documentclass[11pt,a4paper]{article}
\usepackage[margin=24mm]{geometry}
\usepackage{amsmath,amssymb,amsthm,hyperref}
\setlength{\emergencystretch}{2em}
\newtheorem{theorem}{Theorem}\newtheorem{lemma}[theorem]{Lemma}
\newcommand{\gen}[1]{\langle#1\rangle}\newcommand{\Om}{\Omega}
\newcommand{\SO}{\mathrm{SO}}\newcommand{\F}{\mathbf F}
\title{Mixed involutions in orthogonal strings of dimension five}
\author{Human Compression project: F001 working proof}
\date{Unreviewed working manuscript}
\begin{document}\maketitle
This draft completes a proposed route from the accepted bound for negative
reflections to a bound for arbitrary involution types. It has not received
a complete mathematical review. In particular no principal F001 recovery
is inferred from its presence in the working directory.

Let $V$ be a nondegenerate orthogonal five-space in odd characteristic,
with symmetric form $B$. A \emph{negative reflection} is $-r_v$, where
$r_v(x)=x-2B(x,v)v/B(v,v)$ and $B(v,v)\ne0$. Its negative eigenspace
has dimension four. Every other nonidentity involution in $\SO(V)$
has a two-dimensional negative eigenspace; we call it a halfturn.
Throughout the endpoint argument we extend scalars to an algebraic closure.

\section{The end pairs}
\begin{lemma}\label{ends}
Suppose $s_0,\ldots,s_{r-1}$, $r\ge6$, is a connected string of
involutions in $\SO(V)$ generating an absolutely irreducible group.
Then $s_0,s_1,s_{r-2},s_{r-1}$ are negative reflections.
\end{lemma}
\begin{proof}
Put $A=\gen{a,b}=\gen{s_0,s_1}$ and
$D=\gen{s_{r-2},s_{r-1}}$. Both are nonabelian dihedral groups and
commute elementwise. Every mark lies in $A,D,C(A)$ or $C(D)$.
Thus a proper subspace preserved by these four groups contradicts
irreducibility. Also if $C(A)$ and $C(D)$ commute elementwise, then
$s_2\in C(D)$ commutes with $s_3\in C(A)$, another contradiction.
We show that $a,b$ must be negative reflections. Reversing the string
then gives the other two marks.

Set $h=ab$. It has determinant one, $h^2\ne1$, and $aha=h^{-1}$.
Its orthogonal primary spaces pair only reciprocal eigenvalues.
Consider first semisimple $h$. If its nonunit reciprocal eigenvalues
are all simple and occupy four dimensions, its orthogonal centralizer
is a torus times signs and is abelian, which cannot contain $D$.
If a nonunit reciprocal pair has multiplicity two, let $E,E'$ be its
two eigenspaces. Choose a basis of $E$ and its image under $a$ as
basis of $E'$. On $E\oplus E'$ the involution $a$ exchanges the
planes and the form has matrix
\[
                 \begin{pmatrix}0&S\\S&0\end{pmatrix},
                 \qquad S=S^{\mathsf T}\text{ invertible}.
\]
The remaining line is fixed by $a$, by determinant one. Thus
$C_{\SO(V)}(A)$ consists of $\operatorname{diag}(T,T,1)$ with
$T^{\mathsf T}ST=S$. The noncommuting involutions of $D$ act on this
two-dimensional multiplicity space as noncommuting orthogonal
reflections; their action is irreducible. The common fixed space of $D$ is
therefore exactly the remaining line, since its action on the four-space
is two copies of that irreducible two-space. Its centralizer preserves
this line, which is already fixed by $A,D,C(A)$. This contradicts
irreducibility.

The remaining semisimple case has one simple nonunit reciprocal pair
on a plane $X$ and a three-dimensional unit-primary complement.
On $X$, $a$ is a reflection, so it has determinant $-1$ there.
If $h$ is the identity on the complement, $a$ has either one or three
negative eigenvalues on that complement. With three, $a,b$ are the
desired negative reflections. With one, the complement splits into
a positive plane $Y$ and a negative line. The full centralizer of
$A$ preserves $Y$ and acts on it as $\mathrm O(Y)$, while acting by
signs on the other constituents. Hence a nonabelian $D$ acts
irreducibly on $Y$ and by scalar characters elsewhere. Its centralizer
also preserves $Y$. This is a common proper invariant plane.
If instead $h$ is negative on a plane $Y$ and positive on a line,
commuting with $h$ already separates these constituents. For $C(A)$
to contain a nonabelian $D$, $a$ must be scalar on $Y$; otherwise
all its joint eigenspaces there are lines and the centralizer is
abelian. Again $D$ acts irreducibly on $Y$ and its centralizer preserves
that plane, giving the same contradiction. This exhausts semisimple
$h$, since a spectrum consisting only of $1,-1$ would give $h^2=1$.

For completeness we now treat every nonsemisimple primary possibility.
An even Jordan block at eigenvalue $1$ or $-1$ occurs with even
multiplicity in an orthogonal space of odd characteristic. One sees
this directly on the leading pairings between equal-sized Jordan
chains: the induced nondegenerate form on their multiplicity space
is alternating when the chain length is even. Hence that multiplicity
is even. This usual orthogonal Jordan restriction leaves in dimension
five only the partitions $(5)$, $(3,1,1)$ and $(2,2,1)$ when the
whole spectrum is positive unit. The nontrivial negative-primary
possibilities are $(3,1)$ or $(2,2)$, with a positive residual line;
or a negative scalar plane together with a positive three-chain.
A nonunit reciprocal pair with two size-two blocks has abelian
centralizer: its action is determined on one block by a polynomial
in that Jordan block. A simple nonunit pair together with a positive
three-chain has a torus times the polynomial centralizer of the
three-chain, again abelian. A regular five-chain also has polynomial,
hence abelian, centralizer. None can contain $D$.

Here are explicit centralizers for all the remaining cases. They also
justify that no unipotent mixing of isomorphic constituents has been
silently discarded.

For a positive $(2,2,1)$ partition use a basis $u_1,u_2,v_1,v_2,z$,
with $B(u_i,v_j)=\delta_{ij}$, $B(z,z)=1$, and
\[
 h=1+N,\quad Nv_1=u_2,\quad Nv_2=-u_1,\quad Nu_i=Nz=0.
\]
The relation $aNa=-N$ forces determinant $-1$ on $\operatorname{im}N$: the
induced alternating form on this plane is changed to its negative.
Choose its two $a$-eigenlines and isotropic dual vectors in the matching
nondegenerate eigenspaces of $a$. Their residual perpendicular line
is fixed by determinant one. Thus
$a=\operatorname{diag}(1,-1,1,-1,1)$ in this basis.
Solving commutation with $N,a$, followed by preservation of $B$, gives
exactly the following matrices $M(t,u)$, $t\ne0$:
\[
\begin{aligned}
 Mu_1&=t u_1,&Mu_2&=t^{-1}u_2,&Mv_2&=t v_2,\\
 Mv_1&=-\tfrac12tu^2u_1+t^{-1}v_1+uz,&
 Mz&=-tu u_1+z.
\end{aligned}                                                \tag{1}
\]
For clarity, commutation with $a$ separates $u_2,v_2$ from
$u_1,v_1,z$; commutation with $N$ fixes the two image lines; their
dual pairings then determine the displayed coefficients. The possible
sign on $z$ is positive because the total determinant is one.
Conversely (1) satisfies all the equations. Its multiplication law is
\[
 M(t,u)M(v,w)=M(tv,w+u/v).
\]
Every nonidentity involution has $t=-1$. Two noncommuting involutions
have distinct parameters, and their product is $M(1,w-u)$.
For $c\ne0$, the image of $(M(1,c)-1)^2$ is exactly $\gen{u_1}$,
since it maps $v_1$ to $-c^2u_1$ and kills the other basis vectors.
Thus $C(D)$ preserves that line. It is already preserved by
$A,D,C(A)$, contradicting irreducibility.

For a positive $(3,1,1)$ partition write
$V=W\perp Y$, $W=\gen{e,z,f}$, where
\[
 B(e,f)=1,\quad B(z,z)=1,\quad Nz=e,\quad Nf=-z,\quad Ne=0,
 \qquad h|_W=1+N+N^2/2.
\]
The line $\operatorname{im}N^2$ has an $a$-eigenvalue
$\alpha=\pm1$. Choose $f$ in that eigenspace with $B(e,f)\ne0$.
Its $N$-chain spans a nondegenerate $a$-stable three-space, giving the
stated perpendicular decomposition after normalization. On $e,z,f$
the signs of $a$ are $\alpha,-\alpha,\alpha$, and its involution
$a_Y$ on $Y$ has determinant $-\alpha$. The whole centralizer is
\[
\begin{aligned}
 Me&=c e,&Mz&=c z,\\
 Mf&=c f+y-\tfrac{B(y,y)}{2c}e,&
 Mu&=Tu-cB(y,Tu)e\quad(u\in Y),
\end{aligned}                                               \tag{2}
\]
where $c=\pm1$, $T\in\mathrm O(Y)$, $\det T=c$, $Ta_Y=a_YT$,
and $a_Yy=\alpha y$. Indeed commutation with $N,a$ first fixes
$Me,Mz$ and the allowed $Y$ component of $Mf$; the equations
$B(Mf,Mf)=0$ and $B(Mf,Mu)=0$ determine the remaining terms.
They conversely verify (2).

If $\alpha=-1$, then $a_Y=\beta I$. When $\beta=-1$, both
$a,b$ are negative reflections, as desired. When $\beta=1$, $y=0$;
$C(A)$ preserves $Y$, and a nonabelian $D$ acts irreducibly on that
plane and scalarly on $W$. Therefore $C(D)$ preserves $Y$ too,
contradicting irreducibility. If $\alpha=1$, $a_Y$ is a reflection.
Write $T=\operatorname{diag}(\lambda,\mu)$ in its positive and
negative lines, and normalize the allowed translation by $u=t/c$.
The group law becomes
\[
 (\lambda_1,\mu_1,u_1)(\lambda_2,\mu_2,u_2)
 =(\lambda_1\lambda_2,\mu_1\mu_2,u_1+\mu_1u_2).
                                                               \tag{3}
\]
This is a central $C_2$ times an additive line extended by inversion.
Noncommuting involutions must both have $\mu=-1$ and distinct
translations. Their product has nonzero translation and possibly a
central sign; its square is a nonzero pure translation. By (2), its
squared nilpotent image is $\gen e$. This line is preserved by
$C(D)$ and already by $A,D,C(A)$, a contradiction.

If the negative primary partition is $(2,2)$, with a positive line,
formula (1) applies but cross-primary mixing is forbidden, so $u=0$.
The centralizer is abelian. For negative partition $(3,1)$ and a
positive line, formula (2) applies with a one-dimensional $Y$ inside
the negative primary space and a separate positive line $Z$. Put
$a_Y=\beta$, $a_Z=\gamma$; determinant one gives
$\alpha\beta\gamma=-1$. If $\beta\ne\alpha$, translations vanish
and the centralizer is abelian. Otherwise (2), with the extra sign
on $Z$ fixed by determinant, has exactly the group structure (3).
The same characteristic-line argument applies. Finally, for a positive
three-chain and a negative scalar plane, cross-primary mixing again
forces $y=0$ in (2). If the restriction of $a$ to the plane is a
reflection its centralizer is abelian; if scalar, every nonabelian
$D$ and its centralizer preserve that plane. Thus this last case is
also reducible. All nonsemisimple cases are exhausted, and the lemma
follows.
\end{proof}

\section{The middle marks and the length of the string}
We now assume the intersection property as well. We acquire from the
reviewed rank-contraction paper L037 the following elementary consequence:
in a connected C-triple $(a,b,c)$ with $[a,c]=1$, the element $b$
does not commute with $ac$. Its complete proof is an actual dependency.
We also acquire the reviewed tridiagonal lemma: a connected string of
distinct negative reflections in dimension five has length at most six;
at length six either end five roots form a basis.

Let $X$ and $Y$ be the root planes of the first and last pairs from
Lemma~\ref{ends}. They are perpendicular. If their intersection is
nonzero, it is the radical line of both planes, and $A,D,C(A),C(D)$
all preserve it. Thus $X\cap Y=0$. Their four-dimensional sum has
radical of dimension at most one, so at most one plane is degenerate.
Reverse the string if necessary so that $Y$ is nondegenerate.

The mark $s_2$ commutes with both last negative reflections. It therefore
acts by one scalar sign on $Y$, because their nonorthogonal root lines
have the same eigenvalue under any commuting orthogonal involution.
On $Y^\perp$ it has determinant one, and is a three-dimensional
halfturn: identity would make it commute with $s_1$.
Let $v$ be its positive axis there and let $v_0$ be the first root.
Commutation with $s_0$ gives either $v\perp v_0$ or $v=v_0$.
In the latter case either $s_2=s_0$, or $s_0s_2$ is a scalar sign
on $Y$ and identity on $Y^\perp$, so commutes with $s_1$.
Both alternatives contradict the C-property. Thus $v\perp v_0$.
If $s_2$ is a halfturn on $V$, it acts negatively on $v_0$,
and $s_0s_2$ is a negative reflection. Otherwise $s_2$ already is
one. The replacement preserves every string commutation and preserves
noncommutation with $s_1$ by the acquired C-triple fact.

The other neighboring middle mark has the same property. If its
commuting root plane $X$ is degenerate, a halfturn commuting with
both roots must fix $X$ pointwise: its nondegenerate negative plane
cannot equal the degenerate plane $X$. It must act negatively on the
last root $v_{r-1}$, for otherwise its nondegenerate positive
three-space would equal the degenerate span $X+\gen{v_{r-1}}$.
Multiplying by $s_{r-1}$ again gives a negative reflection.

If $r\ge8$, the adjacent marks $s_3,s_4$ lie in $C(A)\cap C(D)$.
When both root planes are nondegenerate, this is the abelian group of
scalar signs on $X,Y,(X+Y)^\perp$, with the last sign forced positive
by determinant. When one plane is degenerate, the three-chain
calculation above shows that the common centralizer is just the scalar
sign on the other plane. In either case adjacent commutation is a
contradiction. Thus $r\le7$.

At $r=7$, the middle mark $s_3$ lies in that common centralizer.
In the degenerate case it already commutes with $s_2$, impossible.
In the nondegenerate case write $V=X\perp Y\perp Z$. To fail to
commute with $s_2\in C(D)$, $s_3$ must have negative sign on $X$;
to fail to commute with $s_4\in C(A)$, it must have negative sign
on $Y$. Its sign on $Z$ is positive, so it is a negative reflection.
Replacing halfturns $s_2,s_4$ by $s_0s_2,s_4s_6$ now gives a
connected string of seven negative reflections generating the same
group. Noncommutation at the inner bonds is unchanged, because the
multiplying endpoint commutes with both original marks there.
The tridiagonal lemma excludes this. Consequently $r\le6$.

At $r=6$ perform these same replacements, if needed, at positions
$2,3$, using endpoints $0,5$. Denote the resulting negative
reflections by $t_0,\ldots,t_5$ and their roots by $v_0,\ldots,v_5$.
They form a connected string generating the same group; either five
consecutive roots are a basis. No intersection property for the changed
marking is asserted or used.


\section{The reflection subgroup of an original end}
Put $G=\Om_5(q)$, with odd $q\ge5$. We retain its simplicity and
absolute irreducibility, and the classical order formula
\[
 \omega(q)=|G|=\tfrac12q^4(q^2-1)(q^4-1),
 \qquad |\mathrm O_5(q)|=4\omega(q).                    \tag{4}
\]
A C-marking of $G$ is connected: disconnected components would give
commuting nontrivial normal factors. Suppose that its rank is six and
form the negative reflections $t_i=-r_{v_i}$ just obtained.
The original end group $H=\gen{s_0,\ldots,s_4}$ is proper, by the
intersection property. It contains negative reflections with roots
$v_0,v_1,v_2,v_4$, since $t_2$ is either $s_2$ or $s_0s_2$.
If $s_3$ is negative it contains the fifth root $v_3$ as well.
Otherwise $s_3=t_3t_5=r_{v_3}r_{v_5}$, and conjugating $t_2$ by
$s_3$ gives a negative reflection with root
\[
                 u=v_2-cv_3,\qquad
                 c=2B(v_2,v_3)/B(v_3,v_3)\ne0.          \tag{5}
\]
The five roots $v_0,v_1,v_2,u,v_4$ form a basis. Their pairing graph
is connected: $v_0,v_1,v_2$ form a path, $u$ pairs nontrivially with
$v_1$ and $v_4$. Therefore, by the acquired invariant-subspace
argument for connected spanning roots, the ordinary reflection group
$K$ generated by all $H$-conjugates of these reflections is absolutely
irreducible. It is normalized by $H$, and its signed image
\[
 \pi(K)\le H,\qquad \pi(x)=\det(x)x.                  \tag{6}
\]
The kernel of $\pi$ on the full orthogonal group is $\{1,-1\}$.

We need here the full finite reflection classification, rather than
its specialization to five generating reflections. We retain the
following dimension-five consequence of Zalesskii--Serezhkin, in the
form described by Monson~\cite[pp.15--18]{M}. Up to extending scalars
and conjugating, an absolutely irreducible finite group generated by
orthogonal reflections in odd characteristic is one of
\begin{enumerate}
\item a group of orthogonal type over a finite field $\F_{q_0}$ of
that characteristic, containing $\Om_5(q_0)$ and preserving a
nondegenerate form over $\F_{q_0}$;
\item the faithful ordinary modular reflection representation of
$A_5$, $B_5$ or $D_5$;
\item the deleted representation of $S_7$ in characteristic seven.
\end{enumerate}
For $A_5$ the characteristic does not divide six. The exceptional
five-dimensional quotient of $E_6$ in characteristic three is already
of orthogonal type, $\mathrm O_2(5,3)$ in the notation of the notes.
The unitary cases in the full classification which admit an invariant
nondegenerate symmetric form reduce to the listed orthogonal cases;
the notes discuss this immediately after Theorem~4.1. We do not use
that theorem's exclusion of cases requiring more than five reflections.
The complete classification remains an external retained result.

\begin{lemma}\label{classical}
The group $K$ cannot be of orthogonal type.
\end{lemma}
\begin{proof}
Conjugate so that $K$ and a nondegenerate invariant form $B_0$ are
defined over $\F_{q_0}$. The form $B$ is a scalar multiple of $B_0$,
by absolute irreducibility. We first show that the $\F_{q_0}$-linear
span of $K$ is the whole matrix algebra $M_5(\F_{q_0})$.
For a reflection with root $v$ over that field the rank-one operator
\[
       P_v=(1-r_v)/2=vB_0(v,-)/B_0(v,v)
\]
is in this span, which is an algebra. The graph of all reflecting
root lines, joined by nonzero pairings, is connected and spans $V$;
otherwise the spans of its components give a proper invariant space.
Products of the $P_v$ along a path are nonzero scalar multiples of
$vB_0(w,-)$ for its two end roots. Choosing five root lines as a
basis now gives all matrix units in that basis. This proves the claim.

Since $H$ normalizes $K$, it normalizes this matrix algebra. We give
the elementary normalizer argument to avoid an implicit field-automorphism
assumption. If $h$ normalizes it, put $F_{ij}=hE_{ij}h^{-1}$, whose
entries are in $\F_{q_0}$. Choose a nonzero vector $w$ over that
field in the image of $F_{11}$. The vectors $w_i=F_{i1}w$ form a
basis, since $F_{ij}w_k=\delta_{jk}w_i$. Let $A$ have these columns.
Then $A\in\mathrm{GL}_5(q_0)$ and $A^{-1}h$ commutes with every
$E_{ij}$, hence equals a scalar $c$. Thus $h=cA$.
The equations preserving $B_0$ give $c^2\in\F_{q_0}$, by any
nonzero entry of $B_0$; determinant one gives $c^5\in\F_{q_0}$.
Consequently $c=c^5/(c^2)^2\in\F_{q_0}$, and
$H\le\mathrm{SO}_5(q_0)$.

The roots $v_0,v_1,v_2$ are therefore lines defined over $\F_{q_0}$.
If $s_3$ is negative, so is the root $v_3$, and the missing root
line $\gen{v_5}$ is the perpendicular of their four-dimensional span.
If $s_3$ is a halfturn, its negative plane is
$\gen{v_3,v_5}$, and
\[
 \gen{v_5}=\gen{v_3,v_5}\cap\gen{v_0,v_1,v_2}^{\perp}.       \tag{7}
\]
The intersection has dimension one because $v_3$ pairs nontrivially
with $v_2$, whereas $v_5$ pairs with none of the first three roots.
Both spaces on the right are defined over $\F_{q_0}$. In either
case $s_5$ is defined over that field too, so
$G=\gen{H,s_5}\le\mathrm{SO}_5(q_0)$ after this conjugation.

By (6), $\omega(q_0)\le|\pi(K)|\le\omega(q)$: the subgroup
$\Om_5(q_0)$ has trivial intersection with the kernel of $\pi$,
since $-1$ has determinant $-1$. Thus $q_0\le q$. If $q_0<q$,
these are powers of the same odd prime, so $q\ge3q_0$ and (4)
gives $\omega(q)>4\omega(q_0)$, contradicting the inclusion of
$G$ in $\mathrm O_5(q_0)$. If $q_0=q$, (6) shows that
$\pi(K)=G\le H$, contradicting the properness of the end group.
\end{proof}

\section{The finite end groups}
We next determine the possible original end $H$, not merely $K$.
In the $A_5$ case $K\cong S_6$, and in the deleted case
$K\cong S_7$. In either representation the ordinary reflections are
precisely the transpositions. Indeed an involution with $k$ disjoint
transpositions has negative eigenspace dimension $k$ on the deleted
or sum-zero module, as the removed trivial constituents are positive.
A linear normalizer therefore induces an automorphism preserving the
transposition set and its commutation graph. Its vertices are the edges
of the complete graph on six or seven letters, with adjacency when
they meet. Every family of at least four pairwise intersecting edges
is contained in a star: after edges $ab,ac$ are chosen, an edge missing
$a$ must be $bc$, leaving at most that triangle. Thus the maximal
cliques of size at least four are exactly the stars. An automorphism
permutes these stars, hence permutes the letters, and their pairwise
intersections determine every transposition. It follows that the
induced group automorphism is inner. After multiplying an element of
$H$ by the appropriate element of $\pi(K)$, it centralizes the
absolutely irreducible group $K$, so is scalar. A scalar in
$\mathrm{SO}_5$ is the identity, and consequently $H=\pi(K)$
in both cases.

For $B_5$ and $D_5$, the largest normal 2-subgroup of $K$ is the
full or even diagonal sign group. Indeed the quotient is $S_5$,
whose largest normal 2-subgroup is trivial. Its five distinct common
weight spaces are the coordinate axes, so every element of $H$
permutes those axes. An orthogonal monomial matrix has entries $\pm1$
in its nonzero positions. The determinant-one monomial group is
\[
              W=E_5\rtimes S_5,\qquad
 E_5=\{(\epsilon_i)\in\F_2^5:\textstyle\sum\epsilon_i=0\},
                                                               \tag{8}
\]
where a permutation is represented by its signed matrix
$\operatorname{sgn}(\sigma)P_\sigma$. Both $B_5$ and $D_5$ have
signed image equal to $W$. Thus $H=W$ in these cases.

\begin{lemma}\label{finiteends}
If $H\cong S_6$ or $H=W$ above, every one of its five original
marks is a negative reflection.
\end{lemma}
\begin{proof}
For $S_6$ we acquire the reviewed classification of independent
extremal strings: five involutions generating $S_6$ independently and
satisfying the string commutations form a Coxeter path, up to an
automorphism. In the present signed representation a transposition
is a negative reflection, whereas an outer image, a triple
transposition, is a halfturn. Since $s_0,s_1,s_4$ are negative
reflections, only the first possibility occurs.

For $W$ we acquire the reviewed even-sign module lemma and the
extremal independent string result for $S_5$ from the F006 recovery.
Choose a minimal subfamily $L$ whose images generate $S_5$. The
module lemma permits at most one mark outside $L$, and Whiston's
bound gives $|L|\le4$. Since all five marks are independent, exactly
four belong to $L$. Their images form a Coxeter path in $S_5$.
Adjacent images do not commute, so their original positions were
consecutive as well. The omitted mark $x$ is therefore an endpoint.
Reverse the string if needed to put it first.

It commutes with the last three members of $L$, whose images generate
an $S_4$. The centralizer of this $S_4$ in $S_5$ is trivial, so
$x$ belongs to $E_5$. It is fixed by that $S_4$. The only nonzero
even-weight vector constant on its orbits of sizes one and four has
weight four. Thus $x$ is a negative reflection.

Put $H_0=\gen L$. Its kernel over $S_5$ is an $S_5$-submodule of
$E_5$. The acquired module lemma says it is zero or all $E_5$,
as the nonzero constant vector has odd weight. The latter would put
$x$ in $H_0$, contrary to independence. Hence $H_0\cong S_5$.
Each mark in $L$ is consequently the unique lift in $H_0$ of a
transposition. The lifts of permutations of its three fixed coordinates
commute with it. They permute those three axes, so the signs of the
mark on these axes are equal. Its swapped pair contributes determinant
$-1$, and total determinant one forces each fixed sign to be $-1$.
The swapped pair contributes one negative eigenvalue, giving four in
all. Thus these four marks are negative reflections too.
\end{proof}

\section{A deleted symmetric end cannot extend}
We give the remaining end argument in coordinates, including the
quotient by the constant line. Suppose $K$ has the deleted $S_7$
representation in characteristic seven, so $H=\pi(K)\cong S_7$.
In the seven-coordinate permutation space put
\[
 T=\{(x_i):\textstyle\sum_{i=1}^7x_i=0\},\qquad
 V=T/\gen{\boldsymbol1},\qquad \boldsymbol1=(1,\ldots,1).
                                                               \tag{9}
\]
The dot product induces a nondegenerate form on $V$.
The negative eigenspace of an involution with $k$ disjoint
transpositions has dimension $k$ in $V$; its determinant is $(-1)^k$.
After signing, its negative dimension is four for $k=1$, two for
$k=2$, and two for $k=3$. Thus the negative reflections in $H$ are
precisely the signed transpositions.

Ignore the scalar signing when discussing the associated permutations.
Relabel coordinates so that
\[
       s_0=(12),\qquad s_1=(23),\qquad s_4=(45).          \tag{10}
\]
This is possible since the first pair does not commute and the last
mark commutes with both. The element $s_3$ centralizes $S_{\{1,2,3\}}$,
so acts only on $\{4,5,6,7\}$ and does not commute with $(45)$.
It is therefore a single or double transposition there. The condition
that $s_2$ commutes with both $(12)$ and $(45)$ gives
\[
 s_2=(12)^\epsilon(45)^\delta\tau,
 \qquad \epsilon,\delta\in\{0,1\},\quad
 \tau\in\{1,(36),(37),(67)\}.                         \tag{11}
\]
Transitivity of $H=S_7$ forces $\tau=(36)$ or $(37)$.
If $s_3$ is a single transposition, normalize it to $(46)$.
Then $\tau=(36)$ leaves coordinate seven fixed, while $\tau=(37)$
leaves the orbits $\{1,2,3,7\}$ and $\{4,5,6\}$ separate.
Neither is transitive. Thus, after relabeling the last four coordinates,
\[
                    s_3=(46)(57).
\]
Conjugating by $(45)(67)$ preserves this and $s_4$ and exchanges
$(36)$ and $(37)$. Hence we may take $\tau=(36)$ in (11).

The negative reflection $s_5$ commutes with $s_0,\ldots,s_3$,
so its root is a common eigenline of these permutations in $V$.
A common eigenline of $(12),(23)$ must be positive for each: neither
one-dimensional negative eigenspace is preserved by the other.
It therefore has a representative with first three coordinates equal,
say $x_1=x_2=x_3=c$. There is no lift ambiguity for a positive
involution eigenvector: if $(g-1)x=a\boldsymbol1$, applying $g$
once more gives $2a=0$, hence $a=0$. A negative eigenvector modulo
the constant line can be made an exact negative eigenvector by
subtracting a suitable constant vector.

If the root has negative eigenvalue for the permutation $s_2$, choose
that exact negative representative. Its first three coordinates are
then zero, as is $x_6$. When $\delta=0$ all coordinates vanish.
When $\delta=1$ the only remaining possibility is a multiple of
$e_4-e_5$, which is not an eigenline of $(46)(57)$, even modulo
constants. Thus the root is positive for $s_2$, and
\[
                   x_6=c,\qquad \delta=1\Longrightarrow x_4=x_5.
                                                               \tag{12}
\]
If it is also positive for $s_3$, then $x_4=x_6=c$ and
$x_5=x_7=d$. The sum-zero condition is $5c+2d=0$, which in
characteristic seven says $d=c$. Its quotient class is zero.
Finally suppose it is negative for $s_3$, and subtract a constant
to obtain an exact negative eigenvector. This preserves (12).
Since $s_3$ fixes the first three coordinates, $c=0$, so
$x_4=x_6=0$ and $x_5=-x_7$. When $\delta=1$ this is zero;
when $\delta=0$ the sole possible root line is
\[
                         \gen{e_5-e_7}.                 \tag{13}
\]
Its ordinary reflection is the transposition $(57)$, so its negative
reflection already belongs to $H$. This contradicts the independence
of $s_5$ from the original end. No such end can extend to a rank-six
marking of $G$.

\begin{theorem}\label{main}
For every odd $q\ge5$, a string C-group marking of $\Om_5(q)$ has
rank at most five.
\end{theorem}
\begin{proof}
The endpoint and middle arguments exclude rank greater than six.
Suppose that rank six occurs. If every mark is a negative reflection,
the acquired reflection bound gives a contradiction. Otherwise at
least one middle mark is a halfturn, since all four end marks are
negative reflections. The original end $H=\gen{s_0,\ldots,s_4}$
contains that halfturn. Its normal reflection subgroup $K$ cannot
be classical by Lemma~\ref{classical}. If it is an ordinary finite
alternative, Lemma~\ref{finiteends} makes all the end marks negative,
again a contradiction. Its only remaining possibility is the deleted
$S_7$ case, excluded by (9)--(13). Hence rank six is impossible.
\end{proof}

\paragraph{Inputs, contributions and present status.}
The full finite reflection classification and the classical group facts
are retained external results.
The previously reviewed negative-reflection bound, the C-triple
argument, and the F006 module and extremal permutation results are
acquired proofs. Their complete required interfaces and finite
certificates must accompany any eventual frozen packet. The new
argument treats all orthogonal endpoint spectra, changes the middle
marks only to obtain spanning roots, and studies the normal reflection
subgroup of an unchanged end. The separate review agent contributed
independent working calculations for the nonsemisimple centralizers;
that contribution must be disclosed in review provenance. This draft
has not yet received complete distinct review. Its accompanying exact
arithmetic checks support the centralizer formulas and the deleted
symmetric obstruction; they do not verify the retained classification
or substitute for review of the uniform proof. This file does not itself
confer a principal F001 disposition,
conjecture credit, or a compression measurement.

\begin{thebibliography}{9}
\bibitem{M} B.~Monson, \emph{Constructions of Symmetric Polytopes},
author-hosted course notes, supplied 54-page edition,
\url{https://bmonson.ext.unb.ca/fields/outline.pdf}, Section~4,
pp.15--19.
\bibitem{ZS} A.~E.~Zalesskii and V.~N.~Serezhkin,
\emph{Finite linear groups generated by reflections},
Mathematics of the USSR-Izvestiya \textbf{17} (1981), 477--503.
\end{thebibliography}
\end{document}
